\documentclass[11pt]{article}

\usepackage[utf8]{inputenc}
\usepackage[T1]{fontenc}
\usepackage{lmodern}
\usepackage{geometry}
\usepackage{amsmath,amssymb,amsfonts,amsthm,mathtools}
\usepackage{bm}
\usepackage{enumitem}
\usepackage{microtype}
\usepackage{hyperref}
\usepackage{titlesec}
\usepackage{booktabs}
\usepackage{array}
\usepackage{setspace}
\usepackage{mathrsfs}
\usepackage{physics}

\geometry{margin=1in}
\hypersetup{colorlinks=true, linkcolor=black, urlcolor=blue, citecolor=black}
\setlist[enumerate]{topsep=0.4em,itemsep=0.35em}
\setlist[itemize]{topsep=0.3em,itemsep=0.25em}
\titleformat{\section}{\large\bfseries}{\thesection.}{0.5em}{}
\titleformat{\subsection}{\normalsize\bfseries}{\thesubsection.}{0.5em}{}
\setstretch{1.06}

\newcommand{\Aether}{\AE{}ther}
\newcommand{\Aflow}{\AE{}ther-flow}
\newcommand{\TheoryName}{\Aflow{} Interpretation of Relativity}
\newcommand{\TheoryProgram}{\Aether{} / \Aflow{} framework}
\newcommand{\Sub}{\mathcal{A}}
\newcommand{\BridgeMetric}{\mathfrak{g}}
\newcommand{\Ell}{\mathcal{L}}

\newtheorem{definition}{Definition}
\newtheorem{proposition}{Proposition}
\newtheorem{remark}{Remark}
\newtheorem{corollary}{Corollary}
\newtheorem{theorem}{Theorem}

\title{\textbf{The \TheoryName{}}\\[0.4em]
\large Higher-Derivative Quartic Compensator Branch Post-Restricted-Verdict Derivational Next-Step and Criteria Note}
\author{}
\date{}

\begin{document}

\maketitle

\begin{abstract}
The tuned exact-square compensator sequence has now reached a sharper endpoint than another same-branch asymptotic repair. At flat-background quadratic scope, complete auxiliary elimination removes the full observer correction on the tuned branch while preserving the bridge metric, the single-metric matter route, and the gapped/auxiliary nonmetric sectors. At nonlinear reduced-system scope, the same tuned branch has an explicit unit-lapse spatial-harmonic \(3+1\) architecture, a first local well-posedness and continuation theorem, and a corrected coercive bootstrap. At asymptotic scope, however, the ledger is now closed in a mixed but disciplined way: the renormalized same-branch note remains negative on generic tuned data, while the invariant momentum-suppressed branch closes positively only as a restricted same-class asymptotic theorem.

This note records what that verdict changes. First, it classifies the momentum-suppressed asymptotic theorem as a restricted positive side result rather than the final derivational claim of the tuned branch. Second, it shows that the active burden is no longer another same-branch asymptotic manuscript. The tuned asymptotic bookkeeping has already reached its honest endpoint at current scope: generic tuned data remain closed, and the only positive tuned theorem is restricted to the invariant branch \(\pi_\sigma\equiv 0\). Third, it identifies the tuned exact-square branch as the active viable derivational branch at present scope precisely because it is the only currently recorded branch that both removes the flat-background observer correction completely and already carries an explicit reduced nonlinear architecture.

The remaining burden is therefore the bridge question again, now on the viable tuned branch. Define the reduced tuned bridge correction tensor by
\[
G_{\mu\nu}[g]
=
8\pi G_N\,T_{\mu\nu}^{\mathrm{matter}}
+ \mathcal C_{\mu\nu}^{\mathrm{tc,br}},
\qquad
g_{\mu\nu}=\Xi^\ast\BridgeMetric_{AB}.
\]
The next proof sequence is then exact and finite: prove that the tuned cancellations survive beyond the current flat/local packaged scope; derive the infrared Einsteinian metric equation on that viable branch; decide whether \(\mathcal C_{\mu\nu}^{\mathrm{tc,br}}\) vanishes or is only gauge/constraint exact, or instead leaves a genuine observer-accessible deformation/remainder; and only if such a remainder survives should a replacement bridge/completion candidate be introduced. The present note does not perform that derivation. Its role is narrower: it closes the tuned asymptotic bookkeeping at current scope and fixes the next exact-closure proof chain before any broader redesign is entertained.
\end{abstract}

\section{Introduction}

The tuned exact-square compensator branch has now been carried through enough of the manuscript sequence that its next burden can be stated without ambiguity. The full flat-background reduction note already showed that, on the tuned branch,
\[
\frac{\alpha_Y^2}{M_Y^2}=\frac{\lambda_4}{M_*^2}
\]
removes the \emph{complete} reduced observer correction after auxiliary elimination rather than only an isolated scalar subsection. The later reduced-system and local-well-posedness notes then turned that coefficient-level branch into an explicit nonlinear architecture with a first honest theorem package. The post-coercive asymptotic sequence finally settled the tuned same-branch bookkeeping itself: generic tuned data remain asymptotically closed negatively, while the invariant momentum-suppressed sub-branch admits a restricted positive asymptotic theorem.

That mixed verdict is strong enough to change the research direction again. The next step is not another same-branch asymptotic repair on the tuned exact-square branch. The generic tuned-data obstruction has already been recorded, and the only currently positive tuned asymptotic result is explicitly restricted to \(\pi_\sigma\equiv 0\). The active burden therefore returns to the broader derivational question: whether the tuned cancellations that remove the observer correction at the recorded flat-background scope actually survive beyond the current flat/local packaged regime strongly enough to yield the infrared Einsteinian observer equation on the viable branch itself.

\paragraph{Framework and claim boundary.}
\TheoryProgram{} denotes the broader \Aether{} / \Aflow{} framework built on the claims that the \Aether{} is the underlying four-dimensional substrate of reality, the \Aflow{} is its intrinsic ordered motion, observed three-dimensional space is the local experiential slice of that deeper substrate, S-time is the experienced order of change arising from matter, light, and the \Aflow{}, and observed expansion is the three-dimensional appearance of deeper four-dimensional ordered motion. Within that framework, \TheoryName{} denotes the exact-closure theory in which the effective gravitational dynamics are adopted to be exactly Einsteinian with universal matter coupling. In the active sequence, \emph{adoption} means use of that established relativistic dynamics without claiming substrate derivation, while \emph{derivation} is reserved for a first-principles recovery from explicit substrate variables.

The present note stays within that boundary. It does not claim that the tuned exact-square branch has already derived Einsteinian gravity beyond the current packaged scope. It does not reopen a generic tuned-data asymptotic theorem. It does not redesign the branch again before the next bridge-level derivation is attempted. It records the restricted tuned asymptotic theorem at its correct status and fixes the next derivational proof sequence on the current viable branch.

\section{Post-Restricted Verdict on the Tuned Branch}

\begin{definition}[Current flat/local packaged tuned scope]
\label{def:packaged}
The \emph{current flat/local packaged tuned scope} consists of the following already recorded results on the tuned exact-square compensator branch:
\begin{enumerate}
    \item the full flat-background quadratic reduction with complete auxiliary elimination, on which the reduced observer correction vanishes identically on the tuning
    \begin{equation}
    \frac{\alpha_Y^2}{M_Y^2}=\frac{\lambda_4}{M_*^2};
    \label{eq:tuning}
    \end{equation}
    \item the explicit unit-lapse spatial-harmonic reduced \(3+1\) system with bridge metric
    \begin{equation}
    g_{\mu\nu}=\Xi^\ast\BridgeMetric_{AB}
    \label{eq:obsmetric}
    \end{equation}
    and matter route
    \begin{equation}
    S_{\mathrm{matter}}[\Xi^\ast\BridgeMetric,\psi]
    =
    S_{\mathrm{matter}}[g,\psi];
    \label{eq:matterroute}
    \end{equation}
    \item the first tuned local well-posedness, continuation, and corrected coercive-bootstrap theorem in the repaired weighted/homogeneous class;
    \item the negative renormalized same-branch asymptotic verdict on generic tuned small data;
    \item the positive momentum-suppressed same-branch \(L_t^1\) remainder verdict and the positive restricted asymptotic-closure theorem on the invariant sub-branch
    \begin{equation}
    \pi_\sigma|_{t=0}=0
    \quad\Longrightarrow\quad
    \pi_\sigma\equiv 0.
    \label{eq:restrictedbranch}
    \end{equation}
\end{enumerate}
\end{definition}

\begin{definition}[Restricted positive asymptotic side result]
\label{def:side}
A theorem on a candidate branch is called a \emph{restricted positive asymptotic side result} if it proves asymptotic closure only on an invariant restricted sub-branch and therefore does not reopen the generic-data asymptotic claim on the ambient branch.
\end{definition}

\begin{proposition}[The tuned asymptotic ledger is now settled at current scope]
\label{prop:ledger}
Within the packaged scope of Definition \ref{def:packaged}, the asymptotic status of the tuned exact-square branch is:
\begin{enumerate}
    \item \textbf{generic tuned-data verdict:} the renormalized same-branch route remains closed negatively on generic small data;
    \item \textbf{restricted tuned verdict:} the momentum-suppressed theorem on \eqref{eq:restrictedbranch} is positive;
    \item \textbf{status of the positive theorem:} the positive result of item (2) is only a restricted positive asymptotic side result in the sense of Definition \ref{def:side}, not a generic tuned-data theorem and not a completed derivational claim.
\end{enumerate}
\end{proposition}

\begin{proof}
Item (1) is exactly the verdict of the renormalized same-branch asymptotic test: transporting the frozen scalar momentum and subtracting the induced linear drift does not restore an \(L_t^1\) observer-side remainder on generic tuned data. Item (2) is exactly the verdict of the momentum-suppressed same-branch and restricted asymptotic-closure notes: on the invariant restriction \(\pi_\sigma\equiv 0\), the completion-specific observer-side source vanishes and the same repaired weighted/homogeneous class closes positively around the same exact flat observer target. Item (3) follows because the theorem of item (2) is proved only on the invariant sub-branch \eqref{eq:restrictedbranch}; it therefore does not alter the negative generic-data statement of item (1). This is precisely the situation described in Definition \ref{def:side}.
\end{proof}

\begin{corollary}[No further same-branch asymptotic repair is the active tuned-branch task]
\label{cor:noasymp}
At current disciplined scope, the next substantive tuned-branch manuscript should not be another same-branch asymptotic repair note.
\end{corollary}

\begin{proof}
Proposition \ref{prop:ledger} leaves no unsettled same-branch asymptotic question of the right size. The generic tuned branch is already closed negatively, while the only positive tuned asymptotic theorem already exists and is restricted to the invariant sub-branch \eqref{eq:restrictedbranch}. Another same-branch asymptotic note would therefore either repeat the generic negative verdict or merely refine the already-recorded restricted side theorem, neither of which addresses the remaining derivational burden.
\end{proof}

\begin{remark}
Corollary \ref{cor:noasymp} does not say that the restricted momentum-suppressed theorem is unimportant. It says that its importance is not to replace the broader derivational claim. Its correct role is to close the tuned asymptotic bookkeeping at current scope while the exact-closure burden moves back to the bridge equation on the viable branch.
\end{remark}

\section{The Active Viable Derivational Branch}

\begin{definition}[Active viable derivational branch]
\label{def:viable}
A candidate branch is called an \emph{active viable derivational branch} at current scope if all of the following hold:
\begin{enumerate}
    \item \textbf{bridge and matter preservation:} it preserves the bridge metric \(g_{\mu\nu}=\Xi^\ast\BridgeMetric_{AB}\) and the single-metric matter route \eqref{eq:matterroute};
    \item \textbf{flat-background observer-correction removal:} after complete flat-background quadratic auxiliary elimination, the reduced observer correction vanishes identically on that branch;
    \item \textbf{explicit nonlinear architecture:} it has an explicit reduced nonlinear branch formulation together with at least the first local/coercive theorem package;
    \item \textbf{no broader recorded derivational failure yet:} beyond the current flat/local packaged scope, the branch has not already been shown to leave a nonzero observer-accessible remainder in the intended recovery regime.
\end{enumerate}
\end{definition}

\begin{proposition}[Status of the tuned exact-square branch]
\label{prop:viable}
The tuned exact-square compensator branch is the active viable derivational branch at current scope in the sense of Definition \ref{def:viable}.
\end{proposition}

\begin{proof}
Item (1) of Definition \ref{def:viable} is exactly the recorded bridge and matter structure of the tuned compensator branch. Item (2) is the verdict of the full flat-background consistency/reduction note: on the tuning \eqref{eq:tuning}, the complete reduced observer correction vanishes after auxiliary elimination. Item (3) is supplied by the reduced-system and local-well-posedness notes, which fix the explicit unit-lapse spatial-harmonic branch architecture and the first local/coercive theorem. Finally, item (4) holds because the later tuned notes settled only the same-branch asymptotic bookkeeping. They did not yet carry the branch through the broader bridge-level derivation beyond the current packaged scope. In particular, no note has yet shown that the tuned branch leaves a nonzero observer-accessible bridge remainder once the next nonflat/infrastructure-level derivation is attempted. Therefore the tuned exact-square branch remains viable and active for the derivational program.
\end{proof}

\begin{remark}
Viability is not derivation. Proposition \ref{prop:viable} does not say that the tuned branch already closes onto the exact Einstein equation beyond the current packaged scope. It says only that this is now the branch on which that question should be answered before any replacement redesign is entertained.
\end{remark}

\section{Exact Remaining Derivational Burden on the Viable Branch}

\begin{definition}[Tuned bridge correction tensor beyond the packaged scope]
\label{def:Ctc}
For the tuned exact-square branch, define the \emph{tuned bridge correction tensor} \(\mathcal C_{\mu\nu}^{\mathrm{tc,br}}\) by the reduced observer equation
\begin{equation}
G_{\mu\nu}[g]
=
8\pi G_N\,T_{\mu\nu}^{\mathrm{matter}}
+ \mathcal C_{\mu\nu}^{\mathrm{tc,br}}
\label{eq:tunedbridgeeq}
\end{equation}
in the intended recovery regime after the tuned branch variables and auxiliaries have been reduced beyond the packaged scope of Definition \ref{def:packaged}.
\end{definition}

\begin{remark}
Definition \ref{def:Ctc} is deliberately broader than the already-recorded flat quadratic cancellation. At the packaged flat-background scope, the tuned branch satisfies \(\mathcal C_{\mu\nu}^{\mathrm{tc,br}}=0\). The open question is whether that cancellation survives once the branch is carried into the next derivational regime rather than only the current flat/local packaged one.
\end{remark}

\begin{definition}[Tuned cancellation-survival problem]
\label{def:survival}
The \emph{tuned cancellation-survival problem} is the task of proving that the exact-square tuned cancellations responsible for the flat-background identity \(\mathcal C_{\mu\nu}^{\mathrm{tc,br}}=0\) remain effective after the additional reductions, nonlinear substitutions, and infrared bookkeeping needed beyond the packaged scope of Definition \ref{def:packaged}.
\end{definition}

\begin{definition}[Infrared Einsteinian metric equation on the viable branch]
\label{def:infrared}
The viable tuned branch is said to satisfy an \emph{infrared Einsteinian metric equation} in a stated regime if the reduced observer equation can be written there as
\begin{equation}
G_{\mu\nu}[g]
=
8\pi G_N\,T_{\mu\nu}^{\mathrm{matter}}
+ \mathcal R_{\mu\nu}^{\mathrm{IR}}
\label{eq:infraredeq}
\end{equation}
with \(\mathcal R_{\mu\nu}^{\mathrm{IR}}\) explicitly identified after the tuned reductions, and where the derivational question is whether
\begin{equation}
\mathcal R_{\mu\nu}^{\mathrm{IR}}\equiv 0
\label{eq:irzero}
\end{equation}
or instead \(\mathcal R_{\mu\nu}^{\mathrm{IR}}\) carries a genuine observer-accessible deformation/remainder.
\end{definition}

\begin{definition}[Pure gauge/constraint-exact remainder]
\label{def:pure}
An observer-side remainder tensor \(\mathcal R_{\mu\nu}\) is called \emph{pure gauge/constraint exact} if there exist a one-form \(Y_\mu\) and reduced constraint quantities \(\mathcal H\), \(\mathcal M_\mu\), \(\mathcal A_a\) such that
\begin{equation}
\mathcal R_{\mu\nu}
=
\nabla_{(\mu}Y_{\nu)}
+ \mathcal P_{\mu\nu}\mathcal H
+ \mathcal P_{\mu\nu}{}^\rho\mathcal M_\rho
+ \sum_a \mathcal P_{\mu\nu}^{(a)}\mathcal A_a,
\label{eq:pure}
\end{equation}
where every term on the right vanishes on reduced solutions modulo the diffeomorphism freedom of \(g_{\mu\nu}\).
\end{definition}

\begin{theorem}[Post-restricted-verdict derivational proof sequence]
\label{thm:sequence}
Assume the current packaged tuned scope of Definition \ref{def:packaged} and the viability statement of Proposition \ref{prop:viable}. Then the next exact-closure proof sequence on the tuned branch is:
\begin{enumerate}
    \item \textbf{extend the tuned cancellations beyond the current packaged scope:} solve the tuned cancellation-survival problem of Definition \ref{def:survival};
    \item \textbf{derive the infrared observer equation on the viable branch:} identify the tensor \(\mathcal C_{\mu\nu}^{\mathrm{tc,br}}\) of \eqref{eq:tunedbridgeeq}, equivalently the infrared remainder \(\mathcal R_{\mu\nu}^{\mathrm{IR}}\) of \eqref{eq:infraredeq}, in the intended recovery regime;
    \item \textbf{decide exact closure versus deformation/remainder:} determine whether the tensor found in item (2) vanishes identically, is only pure gauge/constraint exact in the sense of Definition \ref{def:pure}, or instead leaves a genuine observer-accessible remainder;
    \item \textbf{choose a replacement bridge/completion candidate only if item (3) fails positively:} a new redesign becomes the active task only after items (1)--(3) show that the viable tuned branch still leaves a nonzero observer-accessible deformation/remainder.
\end{enumerate}
\end{theorem}

\begin{proof}
Corollary \ref{cor:noasymp} shows that another same-branch asymptotic note is no longer the right next task. Proposition \ref{prop:viable} shows that the tuned exact-square branch is still the branch on which the derivational question should be pursued, because it uniquely combines full flat-background observer-correction removal with an explicit nonlinear architecture. Therefore the remaining burden must be phrased at the level of the reduced observer equation beyond the current packaged scope, namely through the tensor \(\mathcal C_{\mu\nu}^{\mathrm{tc,br}}\) of Definition \ref{def:Ctc}.

Item (1) is forced because the recorded tuned cancellation is presently a flat/local packaged result, while the next question asks whether that cancellation survives in the broader derivational regime actually relevant to recovery. Item (2) is then the first honest bridge-level computation once the survival question is posed: derive the reduced infrared observer equation itself rather than infer it from the asymptotic bookkeeping. Item (3) is exactly the derivational verdict: exact closure requires the observer-side remainder to vanish or to be only pure gauge/constraint exact, while any nonzero observer-accessible remainder leaves a deformation rather than exact recovery. Item (4) follows because a replacement redesign is justified only if the current viable branch has first been carried through the derivational test and shown to fail there. Before that point, replacing the branch would abandon the only presently viable exact-square candidate without settling its actual bridge-level status.
\end{proof}

\begin{corollary}[Replacement-trigger criterion]
\label{cor:trigger}
The restricted momentum-suppressed theorem and the generic tuned-data obstruction do \emph{not} by themselves trigger a replacement bridge/completion candidate. The trigger is narrower: a replacement becomes active only after Theorem \ref{thm:sequence} identifies a nonzero observer-accessible remainder in the reduced infrared observer equation on the viable tuned branch.
\end{corollary}

\begin{proof}
The restricted momentum-suppressed theorem is only a side result by Proposition \ref{prop:ledger}, and the generic tuned-data obstruction is only an asymptotic statement on the unrestricted tuned branch. Neither settles the bridge-level derivational question posed by \eqref{eq:tunedbridgeeq}. Theorem \ref{thm:sequence} therefore makes a replacement candidate downstream of the explicit bridge-level failure verdict, not upstream of it.
\end{proof}

\section{What the Next Manuscript Is Not}

\begin{proposition}[Current non-tasks]
\label{prop:nontasks}
At current scope, none of the following should be treated as the next substantive manuscript burden:
\begin{enumerate}
    \item another same-branch asymptotic repair on the tuned exact-square branch;
    \item a broader generic tuned-data asymptotic theorem;
    \item a replacement bridge/completion note introduced before the bridge-level derivation of Theorem \ref{thm:sequence};
    \item an editorial consolidation pass on the archived derivative-mixing or constraint-assisted branches, or on the standalone exact-closure note.
\end{enumerate}
\end{proposition}

\begin{proof}
Item (1) is ruled out by Corollary \ref{cor:noasymp}. Item (2) is ruled out because Proposition \ref{prop:ledger} keeps the generic tuned-data asymptotic claim closed negatively at current scope. Item (3) is ruled out by Corollary \ref{cor:trigger}: replacement is downstream of an explicit bridge-level remainder verdict. Item (4) does not belong to the derivational proof sequence of Theorem \ref{thm:sequence}; it is a settled tracking/editorial matter rather than the active theory burden.
\end{proof}

\begin{remark}
Proposition \ref{prop:nontasks} is not a claim that those topics are permanently forbidden. It says only that they are not the next disciplined step for the active theory line. The current burden is narrower and more exact: carry the viable tuned branch through the next bridge derivation before reopening redesign or editorial questions.
\end{remark}

\section{Conclusion}

The tuned exact-square branch has now reached a post-restricted-verdict status. The generic tuned asymptotic claim remains closed negatively, while the invariant momentum-suppressed branch supplies a positive theorem only as a restricted side result. That settles the tuned asymptotic bookkeeping at current scope. It does not settle the derivational claim.

The correct next step is therefore exact and branch-specific. The tuned exact-square branch remains the active viable derivational branch because it is the only currently recorded branch that both removes the full flat-background observer correction and already has an explicit reduced nonlinear architecture. The next manuscript should now test whether those tuned cancellations survive beyond the current flat/local packaged scope, derive the infrared Einsteinian observer equation on that viable branch, decide whether exact bridge closure really follows or whether a genuine deformation/remainder survives, and only then, if needed, specify the next replacement bridge/completion candidate.

\end{document}
